\section{\bt{Related Work}{相关工作}}
\label{sec:related}
\paragraph{\bt{Agent memory}{智能体记忆}}
\bi{Agent memories retain facts and conversations~\cite{chhikara2025mem0,packer2023memgpt}, workflows~\cite{wang2025awm}, strategies~\cite{zhao2024expel,ouyang2026reasoningbank}, playbooks~\cite{zhang2026ace}, notes~\cite{xu2025amem}, and plans~\cite{zhang2025apc}. Multi-agent memories share trajectories or intermediate steps~\cite{zhang2025gmemory,fioresi2026learning} and govern access and writes~\cite{rezazadeh2025collaborative,margalit2026governed}; recent systems address invalidation, version checks, transactions, and repair~\cite{wu2026invalidation,chen2026fresh,li2026memtx,zhao2026memorepair}. \system\ makes executable database definitions and their data-dependent validation evidence the objects of sharing, connecting reuse to the database state on which a consumer operates.}
{智能体记忆保存事实与对话~\cite{chhikara2025mem0,packer2023memgpt}、工作流~\cite{wang2025awm}、策略~\cite{zhao2024expel,ouyang2026reasoningbank}、手册~\cite{zhang2026ace}、笔记~\cite{xu2025amem}和计划~\cite{zhang2025apc}。多智能体记忆共享轨迹或中间步骤~\cite{zhang2025gmemory,fioresi2026learning}，管理访问和写入~\cite{rezazadeh2025collaborative,margalit2026governed}；近期系统处理失效、版本检查、事务和修复~\cite{wu2026invalidation,chen2026fresh,li2026memtx,zhao2026memorepair}。\system\ 将可执行数据库定义及其依赖数据的验证证据作为共享对象，把复用与消费者操作的数据库状态联系起来。}

\paragraph{\bt{Database exploration memory}{数据库探索记忆}}
\bi{AgentSM structures prior or synthesized trajectories, retrieves similar questions within a database, and constructs composite tools from recurring tool sequences~\cite{biswal2026agentsm}. It reduces repeated exploration and identifies shared-memory consistency and error propagation as future work. \system\ addresses continued reuse through executable definitions with explicit data conditions: compilation binds the definition to a new question, while shared validation and maintenance establish whether its conditions still hold after updates. This couples reuse of the computation with reuse of the evidence supporting it.}
{AgentSM 将既往或合成轨迹结构化，在同一数据库内检索相似问题，并从重复的工具序列构造复合工具~\cite{biswal2026agentsm}，减少重复探索；它将共享记忆一致性与错误传播列为未来工作。\system\ 通过附带显式数据条件的可执行定义支持持续复用：编译将定义绑定到新问题，共享验证与维护则判定其条件在更新后是否仍成立，从而将计算的复用与支持该计算的证据复用结合起来。}

\paragraph{\bt{Data agents and business semantics}{数据智能体与业务语义}}
\bi{Persistent database knowledge~\cite{liu2026agents,luo2026dataagents} and explicit business semantics~\cite{lei2025spider2,floratou2024nl2sql,li2023bird,sequeda2024kg,huo2026birdinteract,jiang2025siriusbi} support data agents. DataLab extracts metadata and calculation logic from SQL/Python history, retrieves knowledge, generates through a DSL, and updates notebook dependencies after edits~\cite{weng2025datalab}. \system\ couples learned computation with executable data conditions, enabling shared maintenance of the definitions used by subsequent agents.}
{持久数据库知识~\cite{liu2026agents,luo2026dataagents}和明确的业务语义~\cite{lei2025spider2,floratou2024nl2sql,li2023bird,sequeda2024kg,huo2026birdinteract,jiang2025siriusbi}为数据智能体提供支撑。DataLab 从 SQL/Python 历史中提取元数据与计算逻辑，检索知识，通过 DSL 生成代码，并在编辑后更新笔记本依赖~\cite{weng2025datalab}。\system\ 将学得的计算与可执行数据条件结合，对后续智能体使用的定义实施共享维护。}

\bi{MetricFlow and Databricks metric views centralize declared metrics and query generation~\cite{dbt2026metricflow,databricks2026metricviews}. MetricFlow validates configurations, the semantic graph, and SQL executability~\cite{dbt2026metricflowvalidation}; these checks do not establish grain or join predicates on the data. Databricks explicitly leaves \code{rely.at\_most\_one\_match} unvalidated at runtime~\cite{databricks2026metricviews}. Both platforms support data maintenance through caching or materialization (Table~\ref{tab:capabilities}). \system\ additionally learns definitions from successful sessions, derives grain and join conditions, and shares their versioned evidence as the data change. Snapshot-bound execution connects these checks to the query snapshot.}
{MetricFlow 与 Databricks 指标视图集中管理声明的指标和查询生成~\cite{dbt2026metricflow,databricks2026metricviews}。MetricFlow 验证配置、语义图和 SQL 的可执行性~\cite{dbt2026metricflowvalidation}；这些检查不确立数据上的粒度或连接谓词。Databricks 明确不在运行时验证 \code{rely.at\_most\_one\_match}~\cite{databricks2026metricviews}。两个平台都通过缓存或物化支持数据维护（表~\ref{tab:capabilities}）。\system\ 进一步从成功会话学习定义，派生粒度与连接条件，并随数据变化共享带版本的证据；绑定快照的执行机制将这些检查与查询快照关联。}

\bi{\emph{Prediction over relational data.} RelBench poses forecasting tasks over multi-table databases~\cite{robinson2024relbench} and is the testbed of relational deep learning~\cite{fey2024rdl}; LLM agents engineer features for such tasks by writing and validating code or SQL~\cite{hollmann2023caafe,guo2024dsagent}. The database community automates data preparation for learning: FeatPilot searches join paths for the features that most improve a model~\cite{liang2025featpilot}, CtxPipe builds preparation pipelines for one table by reinforcement learning~\cite{gao2024ctxpipe}, DeepPrep explores operator pipelines with an LLM agent over a search tree~\cite{fan2026deepprep}, and KnowTrans transfers knowledge between data-preparation tasks by fine-tuning~\cite{ge2025knowtrans}. \system\ shares executable task programs and feature definitions across sessions on the same database. Programs provide retraining starting points, while definitions support composition across compatible tasks. Compilation applies per-instance temporal bounds, and consumer sessions validate predictive quality and audit temporal leakage.}
{\emph{关系数据上的预测。} RelBench 在多表数据库上提出预测任务~\cite{robinson2024relbench}，也是关系深度学习的实验平台~\cite{fey2024rdl}；LLM 智能体通过编写并验证代码或 SQL 为这类任务构建特征~\cite{hollmann2023caafe,guo2024dsagent}。数据库领域为学习任务自动做数据准备：FeatPilot 搜索对模型提升最大的连接路径特征~\cite{liang2025featpilot}，CtxPipe 用强化学习为单张表构建准备流水线~\cite{gao2024ctxpipe}，DeepPrep 用 LLM 智能体在搜索树上探索算子流水线~\cite{fan2026deepprep}，KnowTrans 通过微调在数据准备任务之间迁移知识~\cite{ge2025knowtrans}。\system\ 在同库会话间共享可执行任务方案与特征定义。任务方案提供重训起点，定义支持兼容任务间的组合。编译添加逐样本时间边界，消费者会话验证预测质量并审计时间泄漏。}

\paragraph{\bt{Sharing work in database systems}{数据库系统中的工作共享}}
\bi{Database work sharing includes multi-query optimization and shared execution~\cite{sellis1988mqo,harizopoulos2005qpipe,giannikis2012shareddb}, result recycling~\cite{ivanova2010recycling}, and caches using semantic descriptions~\cite{dar1996semantic}, update-driven invalidation~\cite{garrod2008ferdinand}, input-dependent keys~\cite{gunda2010nectar}, leases~\cite{nishtala2013memcache}, and transactional consistency~\cite{ports2010txcache}. View maintenance is incremental~\cite{gupta1995maintenance,blakeley1986efficiently}, deferred~\cite{colby1996deferred}, or lazy~\cite{zhou2007lazy}; dependency traces support reusable computations~\cite{mokhov2018build,hammer2014adapton}. \system\ applies these principles to validation evidence derived from learned definitions, sharing condition results across consumers at matching snapshot versions (Lemma~\ref{lem:reuse}).}
{数据库工作共享包括多查询优化与共享执行~\cite{sellis1988mqo,harizopoulos2005qpipe,giannikis2012shareddb}、结果回收~\cite{ivanova2010recycling}，以及采用语义描述~\cite{dar1996semantic}、更新驱动失效~\cite{garrod2008ferdinand}、依赖输入的键~\cite{gunda2010nectar}、租约~\cite{nishtala2013memcache}和事务一致性~\cite{ports2010txcache}的缓存。视图维护可采用增量~\cite{gupta1995maintenance,blakeley1986efficiently}、延迟~\cite{colby1996deferred}或惰性方式~\cite{zhou2007lazy}；依赖轨迹支持计算复用~\cite{mokhov2018build,hammer2014adapton}。\system\ 将这些原则用于从学得定义派生的验证证据，在快照版本匹配时跨使用者共享条件结果（引理~\ref{lem:reuse}）。}

\paragraph{\bt{Constraints and data validation}{约束与数据验证}}
\bi{Summarizability research identifies disjointness and completeness as conditions for correct aggregation~\cite{lenz1997summarizability,mazon2009summarizability,horner2004additivity}, integrity checking specializes verification to the updates that can violate a constraint~\cite{nicolas1982integrity,ceri1990constraint}, and data-validation systems and pipeline tests check declared constraints on tables~\cite{schelter2018deequ,breck2019validation,dbt2026tests}. Building on foundational summarizability and integrity-checking primitives, \system\ automatically derives these conditions from learned definitions and integrates their shared verification into query admission and repair.}
{汇总性研究指出不相交与完整是聚合正确的条件~\cite{lenz1997summarizability,mazon2009summarizability,horner2004additivity}，完整性检查把验证限定在可能违反约束的更新上~\cite{nicolas1982integrity,ceri1990constraint}，数据验证系统与管道测试在表上检查声明的约束~\cite{schelter2018deequ,breck2019validation,dbt2026tests}。基于汇总性与完整性检查的基础机制，\system\ 自动从学得的定义推导这些条件，并将条件的共享验证纳入查询准入和修复。}
